\chapter{Introducción}\label{chapter:introduccion}

Este trabajo fin de máster propone el diseño, desarrollo y evaluación de un sistema
que transforma dictados médicos en informes clínicos completos y normalizados utilizando 
Modelos de Lenguaje de Gran Escala (\englishText{LLMs}). Si bien los \englishText{LLMs} tienen 
un gran potencial para automatizar y estructurar esta información, su integración 
clínica propone tres retos fundamentales: la protección de la privacidad bajo marcos 
legales estrictos, la eliminación de alucinaciones para garantizar la integridad de la información  
y la adopción de estándares internacionales para asegurar la 
interoperatividad. Para ello, se propone un sistema que integra las siguientes etapas:

\boldparagraph{Captura y transmisión segura}{
    Establece un canal de comunicación protegido para el flujo de información sensible desde el punto de dictado hasta el núcleo de procesamiento.
}

\boldparagraph{Transcripción automatizada}{
    Conversión del lenguaje natural hablado a texto mediante modelos especializados de reconocimiento de voz.
}


\boldparagraph{Procesamiento y normalización clínica}{
    Identificación de conceptos médicos y su vinculación con terminologías estándar internacionales (\englishText{SNOMED CT, CIE-10-ES}) para garantizar un lenguaje común.
}

\boldparagraph{Generación de informe con \englishText{LLMs}}{
    Emplea Modelos de Lenguaje (\englishText{LLM}) para sintetizar la 
    información extraída en un informe clínico estructurado y coherente.
}

\boldparagraph{Salida estructurada (\englishText{JSON}/\englishText{FHIR})}{
    Transforma el informe final y las entidades codificadas en formatos de datos 
    legibles por máquinas. Se prioriza el estándar \englishText{HL7 FHIR} (\englishText{Fast Healthcare Interoperability Resources}) para 
    facilitar la integración directa con sistemas de Historia Clínica Electrónica 
    (\englishText{HCE}) y asegurar el intercambio de información entre diferentes 
    niveles asistenciales.
}

\section{Motivación}\label{section:motivacion}

La carga administrativa en el sector sanitario es crítica: 
la redacción de informes consume entre el 25\% y el 40\% del tiempo clínico, 
contribuyendo al agotamiento profesional y reduciendo el tiempo de atención al paciente. 
Aunque el dictado de voz es habitual, su dependencia de la transcripción manual y 
edición posterior genera flujos de trabajo ineficientes.

\section{Objetivos}\label{section:objetivos}

Diseñar, implementar y evaluar un sistema de procesamiento de lenguaje natural (\englishText{Natural Language Processing, NLP}) 
\englishText{end-to-end} para la generación automatizada de informes clínicos estructurados 
a partir de dictados médicos en español. El sistema deberá asegurar la interoperatividad 
mediante el estándar \englishText{HL7 FHIR} (\englishText{Fast Healthcare Interoperability Resources}) y garantizar el cumplimiento estricto 
de las normativas vigentes de privacidad de datos de salud.

\subsection{Objetivos específicos}\label{section:objetivos-especificos}

Para alcanzar el objetivo general, se plantean los siguientes objetivos específicos,
los cuales definen los umbrales de desempeño esperados para el sistema

\begin{itemize}
    \item Lograr una tasa de error de palabra (\englishText{Word Error Rate, WER}) $\le$ 10--12\,\% 
    en la transcripción de dictados médicos en entornos controlados.
    \item Alcanzar una puntuación \englishText{F1} $\ge$ 0.85 en la extracción de entidades clínicas clave 
    (diagnósticos, procedimientos, fármacos, dosis y alergias).
    \item Lograr que al menos el 80\,\% de los informes generados sean aceptados por el personal médico ó un agente de IA 
    con un máximo de dos ediciones menores.
    \item Generar salidas compatibles con el estándar \englishText{FHIR}, asegurando más del 95\,\% de 
    validaciones automáticas exitosas sin errores críticos.
\end{itemize}
